% 原书分章PDF第11—15页：21.3及三种推断方法全文。
\section{推断}

Sofrygin 和 van der Laan（2017）及 van der Laan（2014a）给出的结果表明，对于第21.1节假定的统计模型 $\mathcal M$，上述 TMLE 是渐近线性估计量。这一结果部分依赖于假设A1，即 $\bm W$ 的弱相依性。具体而言，该假设规定，每个 $W_i$ 至多依赖于 $\bm W=(W_1,\ldots,W_N)$ 中另外 $K$ 个观测，其中 $K$ 固定。用 $\sigma_0^2$ 表示该 TMLE 的真实渐近方差；$\sigma_0^2$ 的具体表达式见 Sofrygin 和 van der Laan（2017）。接下来，我们的目标是估计 $\sigma_0^2$，从而开展统计推断。

然而需要注意，尽管所假定的统计模型 $\mathcal M$ 使我们能够构造一致且渐近线性的 TMLE，同一个模型 $\mathcal M$ 却不允许估计其方差，也就是说，在 $\mathcal M$ 下，$\sigma_0^2$ 不可识别。具体而言，除弱相依性以外，该模型不假定我们还知道 $\bm W$ 联合分布的其他信息。但仅凭弱相依性假设，并不能得到 $(W_1,\ldots,W_N)$ 的有效联合似然。因此，我们无法从 $\bm W$ 的联合分布中进行任何形式的重抽样。由此，要有效估计 $\sigma_0^2$，通常需要在原有假设之外，对 $\bm Q_{\bm W}$ 施加额外建模限制。

反过来说，如果能够假定一个限制更强的统计模型，使我们得到 $\bm W$ 的有效似然并能拟合其联合分布，那么也就可以估计方差 $\sigma_0^2$。在这种情况下，可以使用 $\sigma_0^2$ 的经验对应估计量，也可以采用参数自助法（parametric bootstrap）；下面将介绍这两种方法。最后，如果不愿对 $\bm Q_{\bm W}$ 施加额外建模限制，则只能得到 $\sigma_0^2$ 上界的估计。虽然下文提出了一种可能的、临时采用的 $\sigma_0^2$ 上界估计，但其理论验证及更详细的分析仍留待未来研究。

下文先举例说明，哪些对 $\bm W$ 的具体建模限制能够让我们一致估计 TMLE 的方差，例如假定 $\bm W$ 独立同分布。随后介绍一种临时方法，在不对 $\bm Q_{\bm W}$ 施加任何额外建模限制时，它可能提供这一方差上界的估计。最后，介绍以全部观测到的 $\bm W$ 为条件的参数推断；它无需对这些变量的联合分布作任何建模假设。

\subsection{基线协变量受限模型中的推断}

考虑 $\bm W$ 独立同分布这一特殊情形。一种有效的推断方法，是根据 $D_{0,i}$ 的代入估计 $D_{N,i}$，使用 $\sigma_0^2$ 的经验对应估计量。将每个 $E[D_{N,i}(O_i^s)D_{N,j}(O_j^s)]$ 按其经验对应物求值，即可得到如下 $\sigma_0^2$ 的估计量：
\[
 \sigma_N^2=\frac{1}{N}\sum_{i,j}R(i,j)
 \bigl[D_{N,i}(O_i^s)D_{N,j}(\bar O_j)\bigr],
\]
其中
\[
 \begin{aligned}
 D_{N,i}(O_i^s)
 &=D_i(\bar Q_N,\bar g_N,\bar g_N^*,\psi_N^*)(O_i^s)\\
 &=\frac{\bar g_N^*}{\bar g_N}(A_i^s\mid W_i^s)
   \bigl[Y_i-\bar Q_N(A_i^s,W_i^s)\bigr]\\
 &\quad+\Bigl[\int_{\bm a}
 \bar Q_N^*(a_i^s(\bm a,\bm W),w_i^s(\bm W))
 \bm g^*(\bm a\mid\bm W)\\
 &\hspace{7em}{}-\Psi_i(\bar Q_N^*,\bm Q_{\bm W,N})\Bigr]
 \end{aligned}
\]
且
\[
 \begin{aligned}
 \Psi_i(\bar Q_N^*,\bm Q_{\bm W,N})
 =\int_{\bm w}\int_{\bm a}&
 \bar Q_N^*(a_i^s(\bm a,\bm w),w_i^s(\bm w))\\[-.2em]
 &\mathord{}\bm g^*(\bm a\mid\bm w)\,d\bm Q_{\bm W,N}(\bm w).
 \end{aligned}
\]
注意，对于 $i=1,\ldots,N$，可以先从 $N$ 个独立同分布的基线协变量中有放回抽样，再按 $\bm g^*$ 抽取 $N$ 个暴露，从而通过蒙特卡洛积分近似每个 $\Psi_i(\bar Q_N^*,\bm Q_{\bm W,N})$。还需指出，当基线协变量独立同分布时，可以推导出这个特定模型真正的 EIC，van der Laan（2014a）第6.2节已给出了这一推导。但实现其相应的代入渐近方差估计量会带来难以承受的计算量，因此我们建议使用上述 $\sigma_0^2$ 的估计量 $\sigma_N^2$。随后可以构造95\%置信区间（confidence interval，CI）$\psi_N^*\pm1.96\sigma_N/\sqrt N$，它将具有正确的渐近覆盖率。

另一种不要求 $\bm W$ 独立同分布的方法，是假定一个为观测 $i=1,\ldots,N$ 指定特定排序的统计模型。这个排序结合 $\bm W$ 的弱相依性，允许我们为 $(W_1,\ldots,W_N)$ 之间的相依关系假定一个特定统计图，从而定义 $\bm W$ 联合似然的唯一分解。对这些单独的似然因子施加额外建模限制之后，可以得到联合分布 $\bm Q_{\bm W}$ 的单一拟合。这又使我们能够对 $\bm W$ 进行重抽样。因此，利用与上述独立同分布情形相同的代入方法，即可得到每个 $E[D_{0,i}(O_i^s)D_{0,j}(O_j^s)]$ 的估计。

\paragraph{参数自助法。}
另一种方法是采用参数自助法估计 $\sigma_0^2$。注意，这种方法仍然需要对 $\bm W$ 的联合分布施加额外限制，例如独立同分布假设。简而言之，参数自助法从观测数据似然已有的联合拟合中，每次重抽取 $N$ 个观测，随后重新拟合一元最不利子模型的参数 $\epsilon$，并重新计算所得 TMLE。将这一过程重复足够多次，即可得到 TMLE 方差的一致估计。

更详细地说，对以 $b=1,\ldots,M$ 编号的每次自助迭代，首先有放回抽取 $N$ 个协变量 $\bm W^b=(W_1^b,\ldots,W_n^b)$，再用已有的模型拟合 $\hat{\bm g}$ 抽取 $N$ 个暴露 $\bm A^b=(A_1^b,\ldots,A_N^b)$，最后根据已有的结局模型拟合 $\bar Q_N$ 抽取 $N$ 个结局 $\bm Y^b=(Y_1^b,\ldots,Y_n^b)$。这里还假定，已分别把汇总函数 $w_i^s(\cdot)$ 与 $a_i^s(\cdot)$ 应用于 $\bm W^b$ 和 $\bm A^b$，构造出相应的自助随机汇总量 $(W_i^{s,b},A_i^{s,b}:i=1,\ldots,N)$。然后，根据已有拟合 $\bar g_N^*/\bar g_N$，用这一自助样本得到 $i=1,\ldots,N$ 的预测权重。

这些自助预测值作为单位层面的权重，用于执行 TMLE 更新，从而得到初始估计量 $\bar Q_N$ 的自助更新。注意，每次自助迭代中唯一需要的模型拟合步骤，就是 TMLE 模型更新，因此这一过程的计算负担大幅减轻。最后，对这些自助 TMLE 估计值计算经验方差，即得方差估计。第21.6节将用模拟研究评估该方法的有限样本性能，并与上述经验对应方差估计量 $\sigma_N^2$ 的性能进行比较。

\subsection{临时采用的方差上界}

为了避免对 $\bm W$ 的联合分布施加额外建模限制，还可以考虑各种临时方法，得到方差的保守估计。举例来说，可以根据如下 $D_{0,i}$ 的估计，构造 $\sigma_0^2$ 的代入估计：
\[
 \begin{aligned}
 D_{N,i}^c(O_i^s)
 &=\frac{\bar g_N^*}{\bar g_N}(A_i^s\mid W_i^s)
   \bigl[Y_i-\bar Q_N(A_i^s,W_i^s)\bigr]\\
 &\quad+\Bigl[\int_{\bm a}
 \bar Q_N^*(a_i^s(\bm a,\bm W),w_i^s(\bm W))
 \bm g^*(\bm a\mid\bm W)-\bar\Psi_N^*\Bigr],
 \end{aligned}
\]
其中
\[
 \bar\Psi_N^*=\frac{1}{N}\sum_{i=1}^N\Psi_i(\bar Q_N^*,\bm Q_{\bm W,N}).
\]
注意，上述 $D_{0,i}$ 的、针对 $i$ 的估计 $D_{N,i}$ 不再保证均值为零，即不一定得到适当的中心化，从而得到如下保守方差估计：
\[
 \frac{1}{N}\sum_{i,j}R(i,j)
 \bigl[D_{N,i}^c(O_i^s)D_{N,j}^c(\bar O_j)\bigr].
\]
我们的模拟表明，对于 $i=1,\ldots,N$，当朋友数 $|F_i|$ 相对均匀时，这个估计量可以相当准确；当网络的朋友数分布偏斜时，则会变得保守（模拟结果未列出）。最后，在更极端的情形中，例如朋友数服从幂律分布时，上述代入方差估计量会过于保守，以至于无法提供有用信息。我们猜测，随着各个 $\Psi_{i,0}$ 的变异性增大，该估计量会变得越来越保守。对这一估计量的更详细分析留待未来研究。

\subsection{条件目标参数的推断}

诚然，对于许多现实的网络数据生成情境，基线协变量 $\bm W$ 独立同分布的假设可能过于严格。同样，上面提出的其他方法也要求作出特定建模假设，而数据未必支持这些假设。现在提出另一种推断方法：放弃感兴趣的边际参数，直接以观测到的基线协变量 $\bm W$ 为条件进行推断。van der Laan（2014a，第8节）以前已经介绍过这种方法；它得到的 TMLE 与本章介绍的 TMLE 相同。此外，这个 TMLE 达到渐近正态性，且其渐近方差可以识别。渐近方差估计量为
\[
 \sigma_N^2(\bm W)=\frac{1}{N}\sum_{i=1}^N
 \bigl[D_{N,Y_i}^2(O_i^s)\bigr],
\]
其中
\[
 D_{N,Y_i}(O_i^s)=\frac{\bar g_N^*}{\bar g_N}(A_i^s\mid W_i^s)
 \bigl(Y_i-\bar Q_N(A_i^s,W_i^s)\bigr).
\]
注意，这个条件 TMLE 无需对基线协变量的分布建模，因此在对 $\bm W$ 弱得多的一组条件下即可达到渐近正态性，见 van der Laan（2014a，第8节）。所以，在基线协变量弱相依的模型中进行条件推断，是一种很有力的替代方法；当我们愿意接受所得推断的条件性解释时，尤其如此。
